% Generated by paper/scripts/generate_tables.py; do not edit by hand.
\begin{table}[!htbp]
\centering
\scriptsize
\resizebox{\linewidth}{!}{%
\begin{tabular}{@{}L{0.18\linewidth}L{0.24\linewidth}L{0.20\linewidth}L{0.20\linewidth}L{0.28\linewidth}@{}}
\toprule
Candidate & Why it is in the paper & Static screen & Feedback/adaptive screen & Paper status \\
\midrule
Current-system-ish baseline & Stylized comparator, not an empirical claim about the live U.S. system. & rank 33,486; score 0.500; floor 0.419 & bridge 7/7 (0.474); adaptive rank 14,704 & Does not beat the leading synthetic portfolios; kept to show the comparison floor. \\
Balanced static winner & Best point-estimate and lower-bound static screen. & rank 1; score 0.656; max regret 0.107 & bridge 1/7 (0.723); adaptive rank 3,320 & Beats the stylized baseline in static and focused-bridge screens, but gate is do not recommend yet. \\
Robustness challenger & Best cross-profile robustness row and strongest stress-regret challenger. & balanced rank 267; robustness rank 1; floor 0.635 & bridge 6/7 (0.709); adaptive rank 7 & Stronger validation target than the static winner under adaptive screening; gate is calibration gray zone, not recommendation. \\
Adaptive gray-zone leader & Highest all-portfolio adaptive-survival score. & static rank 125; score 0.640; max regret 0.083 & bridge 4/7 (0.713); adaptive rank 1 & Best current validation target, not a recommendation; gate is calibration gray zone, not recommendation. \\
\bottomrule
\end{tabular}
}
\caption{Reader-facing synthesis of the stylized current-system baseline and the main challenger portfolios. ``Beats'' means stronger in the current synthetic screens, not empirically proven superior.}
\end{table}
